\section{USAMO 1993}
        \begin{enumerate}
        	\item (\textit{USAMO 1993}) For each integer $n\ge2$, determine, with proof, which of the two positive real numbers $a$ and $b$ satisfying


 $a^n = a + 1, \quad b^{2n} = b + 3a$

 is larger.


 



\textbf{Solution 1}

First, note that $a, b > 1$, since otherwise $a^n-a$ or $b^{2n}-b$ would be nonpositive. Then, let $f(x)=x^{2n}-x$. Since $f$ is strictly increasing on the interval $(1, \infty)$, we have 
 \begin{align*} (a-1)^2>0 \implies a^2+2a+1>4a \implies a^{2n}>4a \\  \implies a^{2n}-a>3a \implies f(a)> 3a = f(b) \implies a> b && \square \end{align*}


 



\textbf{Solution 2}

Square and rearrange the first equation and also rearrange the second.
 \begin{align} a^{2n}-a&=a^2+a+1\\ b^{2n}-b&=3a \end{align}It is \href{https://artofproblemsolving.com/wiki/index.php?title=Trivial\_Inequality}{trivial} that
 \begin{align*} (a-1)^2 > 0 \tag{3} \end{align*}
since $a-1$ clearly cannot equal $0$ (Otherwise $a^n=1\neq 1+1$).  Thus
 \begin{align*} a^2+a+1&>3a \tag{4}\\ a^{2n}-a&>b^{2n}-b \tag{5} \end{align*}
where we substituted in equations (1) and (2) to achieve (5).  Notice that from $a^{n}=a+1$ we have $a>1$.  Thus, if $b>a$, then $b^{2n-1}-1>a^{2n-1}-1$.  Since $a>1\Rightarrow a^{2n-1}-1>0$, multiplying the two inequalities yields $b^{2n}-b>a^{2n}-a$, a contradiction, so $a> b$.  However, when $n$ equals $0$ or $1$, the first equation becomes meaningless, so we conclude that for each integer $n\ge 2$, we always have $a>b$.

 



\textbf{Solution 3}

Define $f(x)=x^n-x-1$ and $g(x)=x^{2n}-x-3a$. By Descarte's Rule of Signs, both polynomials' only positive roots are $a$ and $b$, respectively. With the Intermediate Value Theorem and the fact that $f(1)=-1$ and $f(2)=2^n-3>0$, we have $a\in(1,2)$.Thus, $-3a\in(-6,-3)$, which means that $g(1)=-3a<0$. Also, we find that $g(a)=a^{2n}-4a$. All that remains to prove is that $g(a)>0$, or $a^{2n}-4a>0$. We can then conclude that $b$ is between $1$ and $a$ from the Intermediate Value Theorem. From the first equation given, $a^{2n}=(a+1)^2=a^2+2a+1$. Subtracting $4a$ gives us $a^2-2a+1>0$, which is clearly true, as $a\neq1$. Therefore, we conclude that $1<b<a<2$.

 


 


	\item (\textit{USAMO 1993}) Let $ABCD$ be a convex quadrilateral such that diagonals $AC$ and $BD$ intersect at right angles, and let $E$ be their intersection. Prove that the reflections of $E$ across $AB$, $BC$, $CD$, $DA$ are concyclic.


 



\textbf{Solution 1}

\textbf{Diagram} <table border="1"><tbody><tr><td>
                    \begin{figure}[H]
                        \centering
                        \includegraphics[width=0.5\textwidth]{Solution/USAMO_I_1993_2_Solution_Figure_1.png}
                    \end{figure}
                    </td></tr></tbody></table>

 


 \textbf{Work} Let $X$, $Y$, $Z$, $W$ be the foot of the altitude from point $E$ of $\triangle AEB$, $\triangle BEC$, $\triangle CED$, 

 


	\item (\textit{USAMO 1993}) Consider functions $f : [0, 1] \rightarrow \mathbb{R}$ which satisfy


 <table><tbody><tr>
<td>     </td><td>(i)</td><td>$f(x)\ge0$ for all $x$ in $[0, 1]$,</td></tr>
<tr>
<td>     </td><td>(ii)</td><td>$f(1) = 1$,</td></tr>
<tr>
<td>     </td><td>(iii)     </td><td>$f(x) + f(y) \le f(x + y)$ whenever $x$, $y$, and $x + y$ are all in $[0, 1]$.</td></tr></tbody></table>Find, with proof, the smallest constant $c$ such that


 $f(x) \le cx$

 for every function $f$ satisfying (i)-(iii) and every $x$ in $[0, 1]$.


 



\textbf{Solution}

My claim: $c\ge2$

 
\textbf{Lemma 1}) $f\left(\left(\frac{1}{2}\right)^n\right)\le\left(\frac{1}{2}\right)^n$ for $n\in \mathbb{Z}, n\ge0$

 For $n=0$, $f(1)=1$ (ii)

 Assume that it is true for $n-1$, then $f\left(\left(\frac{1}{2}\right)^{n}\right)+f\left(\left(\frac{1}{2}\right)^{n}\right)\le f\left(\left(\frac{1}{2}\right)^{n-1}\right)\le \left(\frac{1}{2}\right)^{n-1}$

 $f\left(\left(\frac{1}{2}\right)^{n}\right)\le \left(\frac{1}{2}\right)^{n}$

 By principle of induction, \textbf{lemma 1 is proven}.

 
\textbf{Lemma 2}) For any $x$, $\left(\frac{1}{2}\right)^{n+1}<x\le\left(\frac{1}{2}\right)^n\le1$ and $n\in \mathbb{Z}$, $f(x)\le\left(\frac{1}{2}\right)^n$.

 $f(x)+f\left(\left(\frac{1}{2}\right)^n-x\right)\le f\left(\left(\frac{1}{2}\right)^{n}\right)\le \left(\frac{1}{2}\right)^{n}$ (lemma 1 and (iii) )

 $f(x)\le\left(\frac{1}{2}\right)^n$ (because $f\left(\left(\frac{1}{2}\right)^n-x\right)\ge0$ (i) )

 


 $\forall 0\le x\le 1$, $\left(\frac{1}{2}\right)^{n-1}\ge2x\ge \left(\frac{1}{2}\right)^n\ge f(x)$. Thus, $c=2$ works.

 


 Let's look at a function $g(x)=\left\{\begin{array}{ll}0&0\le x\le \frac{1}{2};\\1&\frac{1}{2}<x\le1;\\\end{array}\right\}$

 It clearly have property (i) and (ii). For $0\le x\le\frac{1}{2}$ and WLOG let $x\le y$, $f(x)+f(y)=0+f(y)\le f(y)$

 For $\frac{1}{2}< x\le 1$, $x+y>1$. Thus, property (iii) holds too. Thus $g(x)$ is one of the legit function.

 


 $\lim_{x\rightarrow\frac{1}{2}^+} cx \ge \lim_{x\rightarrow\frac{1}{2}^+} g(x)=1$

 $\frac{1}{2}c>1$

 $c>2$ but approach to $2$ when $x$ is extremely close to $\frac{1}{2}$ from the right side.

 $\mathbb{Q.E.D}$

 


	\item (\textit{USAMO 1993}) Let $a$, $b$ be odd positive integers. Define the sequence $(f_n)$ by putting $f_1 = a$,
$f_2 = b$, and by letting $f_n$ for $n\ge3$ be the greatest odd divisor of $f_{n-1} + f_{n-2}$.
Show that $f_n$ is constant for $n$ sufficiently large and determine the eventual
value as a function of $a$ and $b$.


 



\textbf{Solution}

\textbf{Part 1}) Prove that $f_n$ is constant for sufficiently large $n$.

 Note that if there is some $f_n=f_{n-1}$ for any $n$, then $\frac{f_{n}+f_{n-1}}{2}=f_n$, which is odd. Thus, $f_{n+1}=f_n=f_{n-1}$ and by induction, all $f_p$ is constant for $p\ge n$.

 Also note that $f_n>0$ since average of $2$ positive number is always positive.

 
Thus, assume for contradiction, $\nexists n$, $f_n=f_{n-1}$.

 Then, $f_{n+1}\le\frac{f_{n}+f_{n-1}}{2}< \max (f_{n},f_{n-1})$, $f_{n+2}\le\frac{f_{n+1}+f_{n}}{2}< \max (f_{n},f_{n+1})\le\max (f_{n},f_{n-1})$

 Thus, $\max (f_{n},f_{n-1})> \max (f_{n+1},f_{n+2})$ and that means that $\max (f_{2n},f_{2n+1})$ is a strictly decreasing function and it must reach $0$ as $n\rightarrow\infty$, which contradict with the fact that $f_n>0$.

 
\textbf{Part 1} proven.

 


 


 \textbf{Part 2}) Show that the constant is $\gcd(a,b)$.

 For any $f_n$ where $\gcd(a,b)=d\ne1$. $f_n=dg_n$ for $g_n$ with the same property except with $g_1=\frac{a}{d}$ and $g_2=\frac{b}{d}$.

 Therefore, if I prove that the constant for any $f_n$ with relatively prime $a$, $b$ is $1$, then I have shown that \textbf{part 2} is true.

 
\textbf{Lemma}) If $\gcd(f_n,f_{n-1})=1$, then $\gcd(f_n,f_{n+1})=1$.

 Assume for contradiction that $\gcd(f_n,f_{n+1})=d\ne1$, since both $f_n$ and $f_{n+1}$ are odd, $d$ is not divisible by $2$.

 $f_{n+1}=\frac{f_n+f_{n-1}}{2^n}$ for some $n\in \mathbb{Z}^+$ such that $f_{n+1}$ is odd.

 $(2^n)f_{n+1}-f_n=f_{n-1}$

 $d(p+q)=f_{n-1}$, where $p$ and $q$ is another integer.

 Thus, $f_{n-1}$ is divisible by $d$ which contradicts with the assumption that $\gcd(f_n,f_{n-1})=1$.

 \textbf{Lemma proven}

 
 By induction, $\gcd(f_n,f_{n-1})=1$ since $\gcd(a,b)=\gcd(f_1,f_2)=1$.

 Since there must exist some $n$ where $f_n=f_{n+1}$ (part 1), $\gcd(f_n,f_{n+1})=f_n=1$.

 $\mathbb{Q.E.D}$

 


	\item (\textit{USAMO 1993}) Let $a_0, a_1, a_2,\cdots$ be a sequence of positive real numbers satisfying $a_{i-1}a_{i+1}\le a^2_i$
for $i = 1, 2, 3,\cdots$ . (Such a sequence is said to be \textit{log concave}.) Show that for
each $n > 1$,


 $\frac{a_0+\cdots+a_n}{n+1}\cdot\frac{a_1+\cdots+a_{n-1}}{n-1}\ge\frac{a_0+\cdots+a_{n-1}}{n}\cdot\frac{a_1+\cdots+a_{n}}{n}$.

 



\textbf{Solution}

Notice that because
 \[(a_0 + a_1 + \dots + a_{n-1})(a_1 + a_2 + \dots + a_n) = (a_0 + (a_1 + \dots + a_{n-1})(-a_0 + (a_0 + a_1 + \dots + a_n))\]
 \[= -a_0^2 + a_0 ((a_0 + a_1 + \dots + a_n) - (a_1 + a_2 + \dots + a_{n-1})) + (a_1 + a_2 + \dots + a_{n-1})(a_0 + a_1 + \dots + a_n)\]
 \[= a_0 a_n + (a_1 + a_2 + \dots + a_{n-1})(a_0 + a_1 + \dots + a_n),\]we may subtract $\dfrac{(a_1 + a_2 + \dots + a_{n-1})(a_0 + a_1 + \dots + a_n)}{n^2}$ from both sides of the inequality and observe that it is sufficient to prove that
 \[\frac{(a_1 + a_2 + \dots + a_{n-1})(a_0 + a_1 + \dots + a_n)}{(n^2 - 1)n^2} \ge \frac{a_0 a_n}{n^2},\]or
 \[(a_1 + a_2 + \dots + a_{n-1})(a_0 + a_1 + \dots + a_n) \ge (n^2 - 1) a_0 a_n.\]

 Fortunately, this is an easy inequality. Indeed, from AM-GM applied on each group of terms we have
 \[(a_1 + a_2 + \dots + a_{n-1})(a_0 + a_1 + \dots + a_n) \ge (n^2 - 1) \sqrt[n-1]{a_1 a_2 \dots a_{n-1}} \sqrt[n+1]{a_0 a_1 \dots a_n},\]and so it suffices to prove
 \[\sqrt[n-1]{a_1 a_2 \dots a_{n-1}} \sqrt[n+1]{a_0 a_1 \dots a_n} \ge a_0 a_n,\]or, after taking both sides to the $(n^2 - 1)$ power, simplifying, and taking the $n$-th root of both sides, to prove
 \[a_1^2 a_2^2 a_3^2 \dots a_{n-1}^2 \ge a_0^{n-1} a_n^{n-1}.\]This easily follows from the Fact that $a_0 a_n \le a_i a_{n-i}$ for $1 \le i \le n-1$. Indeed, we are given that
 \[a_0 a_2 \le a_1^2\]
 \[a_1 a_3 \le a_2^2\]
 \[a_2 a_4 \le a_3^2\]
 \[\dots\]
 \[a_{n-3} a_{n-1} \le a_{n-2}^2\]
 \[a_{n-2} a_n \le a_{n-1}^2.\]Multiply all inequalities together and cancel $a_1, a_2^2, a_3^2, \dots, a_{n-2}^2, a_{n-1}$ to give $a_0 a_n \le a_1 a_{n-1}$. Similarly, by multiplying all inequalities except the first and the last, we deduce that $a_2 a_{n-2} \ge a_1 a_{n-1} \ge a_0 a_n$, and a simple induction argument proves the verity of the Fact for $1 \le i \le \frac{n}{2}$, and so by the Commutative Property the Fact is true for all $1 \le i \le n-1$, as desired. Now multiply each inequality of the Fact for $i = 1, 2, 3, \dots, n-1$ to give the desired result.

 Note: The Fact can be generalized into a Lemma: $a_x a_w \le a_y a_z$ whenever $x \le y \le z \le w$ and $x + w = y + z$. The proof is similar to that of the Fact and is left as an exercise to the reader.

 --User suli, April 15, 2015.

 


\end{enumerate}